% !PRINTED 470
\weberpage{470}

%<<CONT
\[
\varphi(\varrho) = \varphi(\varrho_a)
\]
ist, so sagen wir, \gesperrt{$\omega$ erlaubt oder gestattet die Substitution $(\varrho, \varrho_a)$ oder $\sigma_a$}.

Die Grössen $\omega$, die ausser der identischen Substitution keine Substitution $\sigma_a$ gestatten, sind die primitiven Elemente des Körpers $\Omega(\varrho)$.

\gesperrt{4. Eine Grösse, die alle Substitutionen $\sigma_a$ gestattet, ist nothwendig ein Element von $\Omega$.}

Denn eine solche Grösse ist mit allen ihren conjugirten Grössen identisch, und genügt also nach \S.\,144 einer Gleichung ersten Grades in $\Omega$.

\begin{center}
\S.\ 147.
\end{center}

\begin{center}
\gesperrt{Zusammensetzung der Substitutionen.}
\end{center}

Wenn wir in irgend einer Function von $\varrho$ die Wurzel $\varrho$ zuerst durch $\varrho_a$ und dann $\varrho_a$ durch $\varrho_b$ ersetzen, so ist der Erfolg derselbe, als ob wir gleich von vornherein $\varrho$ mit $\varrho_b$ vertauscht hätten. Setzen wir also
\weq{1}{\sigma = (\varrho, \varrho_a), \quad \sigma' = (\varrho_a, \varrho_b), \quad \sigma'' = (\varrho, \varrho_b),}
so ist es für das Ergebniss gleichgültig, ob wir in allen Zahlen des Körpers $\Omega(\varrho)$ zuerst die Substitution $\sigma$ und dann die Substitution $\sigma'$ ausführen, oder ob wir für einmal die Substitution $\sigma''$ machen.

\gesperrt{1. Wir nennen daher $\sigma''$ aus $\sigma$ und $\sigma'$ componirt oder zusammengesetzt, und bezeichnen diese Beziehung durch die symbolische Gleichung}
\weq{2}{\sigma \sigma' = \sigma''.}
Da wir nach dem Satz 3.\ des vorigen Paragraphen in der Bezeichnung einer Substitution $(\varrho_h, \varrho_k)$ das erste oder das zweite Element beliebig wählen können, so lassen sich $\varrho_a, \varrho_b$ so auswählen, dass $\sigma, \sigma'$ zwei beliebig gegebene unter den $\mu$ Substitutionen des Körpers $\Omega(\varrho)$ sind; $\sigma''$ ist aber dadurch völlig bestimmt. Ebenso ist aber auch, wenn $\sigma$ und $\sigma''$ gegeben sind, $\sigma'$ eindeutig bestimmt. Denn wählen wir $\varrho$ beliebig, so ist $\varrho_a$ durch $\sigma$ völlig bestimmt und $\varrho_b$ durch $\sigma''$, womit auch $\sigma' = (\varrho_a, \varrho_b)$ gegeben ist. Ist endlich $\sigma'$ und $\sigma''$ gegeben, so ist $\sigma$ eindeutig bestimmt, da zunächst $\varrho_b$ durch $\sigma''$ und dann $\varrho_a$ durch $\sigma'$ bestimmt ist. Wir haben daher:
%CONT>>
